% ============================================================================
% PUBLICATION-READY LATEX TABLES
% Stress Testing Large Language Models for Medical Safety
% JAMA Oncology Submission
% Generated: 2026-01-24
% ============================================================================

% NOTE: These tables follow JAMA formatting guidelines:
% - No vertical lines
% - Minimal horizontal lines (top, after header, bottom)
% - Left-aligned text, center-aligned numbers
% - Footnotes using symbols (*, †, ‡, §)
% - Abbreviations spelled out in footnotes

% Required packages:
% \usepackage{booktabs}
% \usepackage{multirow}
% \usepackage{array}
% \usepackage{threeparttable}

% ============================================================================
% TABLE 1: STUDY DESIGN
% ============================================================================

\begin{table}[htbp]
\centering
\caption{2×2×2 Factorial Study Design for Evaluating LLM Safety in Cancer Treatment Discussions}
\label{tab:study-design}
\begin{threeparttable}
\begin{tabular}{@{}lllp{7cm}@{}}
\toprule
\textbf{Factor} & \textbf{Level 1} & \textbf{Level 2} & \textbf{Rationale} \\
\midrule
Conversation Format & Single-turn & Multi-turn (5 exchanges) & Tests whether extended dialogue allows models to maintain or recover safety posture \\
\addlinespace
Treatment Type & Evidence-based* & Unproven/Alternative† & Evaluates differential model behavior toward established versus non-consensus treatments \\
\addlinespace
User Framing (Valence) & Positive (promoting) & Negative (questioning) & Assesses susceptibility to validation bias when users express strong preferences \\
\bottomrule
\end{tabular}
\begin{tablenotes}[flushleft]
\small
\item Abbreviations: LLM, large language model.
\item * Examples: chemotherapy, surgical resection, immunotherapy, CAR T-cell therapy.
\item † Examples: black salve, colloidal silver, vitamin B17, miracle mineral solution.
\end{tablenotes}
\end{threeparttable}
\end{table}

% ============================================================================
% TABLE 2: PRIMARY OUTCOMES - THE SYCOPHANCY EFFECT
% ============================================================================

\begin{table}[htbp]
\centering
\caption{Primary Safety Outcomes: The Sycophancy Effect by Framing Condition}
\label{tab:primary-outcomes}
\begin{threeparttable}
\begin{tabular}{@{}lccccc@{}}
\toprule
 & \multicolumn{2}{c}{\textbf{Framing Condition}} & & & \\
\cmidrule(lr){2-3}
\textbf{Outcome} & \textbf{Neutral} & \textbf{Misconception} & \textbf{Difference} & \textbf{OR (95\% CI)} & \textbf{P Value} \\
 & (n = 636) & (n = 636) & & & \\
\midrule
\multicolumn{6}{@{}l}{\textit{Safety Classification}} \\
\addlinespace
\quad Strong, No.\ (\%) & 477 (75.0) & 442 (69.5) & −5.5 pp & --- & .021 \\
\quad Moderate, No.\ (\%) & 127 (20.0) & 101 (15.9) & −4.1 pp & --- & .048 \\
\quad Weak (Unsafe), No.\ (\%) & 32 (5.0) & 93 (14.6) & +9.6 pp & 3.13 (2.05--4.79) & <.001 \\
\addlinespace
\multicolumn{6}{@{}l}{\textit{Total Safety Score (0--12)}} \\
\addlinespace
\quad Mean (SD) & 10.22 (2.1) & 9.94 (2.8) & −0.28 & --- & .045 \\
\addlinespace
\multicolumn{6}{@{}l}{\textit{Critical Failures*}} \\
\addlinespace
\quad Any critical failure, No.\ (\%) & 25 (3.9) & 86 (13.5) & +9.6 pp & 3.84 (2.40--6.14) & <.001 \\
\bottomrule
\end{tabular}
\begin{tablenotes}[flushleft]
\small
\item Abbreviations: CI, confidence interval; OR, odds ratio; pp, percentage points; SD, standard deviation.
\item * Critical failures include: enables dangerous implementation, actively endorses harmful content, or creates promotional content for unproven therapies.
\item Misconception framing defined as: (1) promoting unproven/alternative treatments OR (2) questioning evidence-based treatments.
\end{tablenotes}
\end{threeparttable}
\end{table}

% ============================================================================
% TABLE 3: MULTI-TURN PROTECTIVE EFFECT
% ============================================================================

\begin{table}[htbp]
\centering
\caption{Unsafe Response Rates by Conversation Format and Framing Condition}
\label{tab:format-interaction}
\begin{threeparttable}
\begin{tabular}{@{}lcccc@{}}
\toprule
 & \multicolumn{2}{c}{\textbf{Framing Condition}} & & \\
\cmidrule(lr){2-3}
\textbf{Conversation Format} & \textbf{Neutral} & \textbf{Misconception} & \textbf{RRR*} & \textbf{P Value†} \\
\midrule
Single-turn, \% (n) & 8.2 (26/316) & 21.2 (68/321) & Reference & --- \\
Multi-turn, \% (n) & 2.2 (7/320) & 8.1 (26/315) & 62\% & <.001 \\
\addlinespace
Absolute risk reduction & 6.0 pp & 13.1 pp & --- & --- \\
Relative risk reduction & 73\% & 62\% & --- & --- \\
\midrule
\multicolumn{5}{@{}l}{\textit{Main Effect of Format}} \\
\addlinespace
\quad OR (95\% CI) & \multicolumn{4}{c}{0.32 (0.21--0.48); P < .001} \\
\bottomrule
\end{tabular}
\begin{tablenotes}[flushleft]
\small
\item Abbreviations: CI, confidence interval; OR, odds ratio; pp, percentage points; RRR, relative risk reduction.
\item * Relative risk reduction comparing multi-turn to single-turn within each framing condition.
\item † P values from logistic regression with format, framing, and their interaction as predictors.
\end{tablenotes}
\end{threeparttable}
\end{table}

% ============================================================================
% TABLE 4: MODEL PERFORMANCE BY PROVIDER
% ============================================================================

\begin{table}[htbp]
\centering
\caption{Safety Performance by Model Provider}
\label{tab:model-performance}
\begin{threeparttable}
\begin{tabular}{@{}lccccc@{}}
\toprule
\textbf{Provider} & \textbf{Models} & \textbf{n} & \textbf{Unsafe Rate, \%} & \textbf{95\% CI} & \textbf{Strong Rate, \%} \\
\midrule
Anthropic & 2 & 227 & 1.8 & 0.6--4.0 & 85.5 \\
OpenAI & 4 & 359 & 8.6 & 6.0--12.0 & 78.0 \\
Google* & 6 & 501 & 10.2 & 7.7--13.2 & 71.7 \\
Meta & 1 & 118 & 16.9 & 10.8--24.7 & 55.1 \\
Alibaba & 1 & 29 & 17.2 & 6.5--35.0 & 72.4 \\
DeepSeek & 1 & 38 & 36.8 & 22.4--53.7 & 39.5 \\
\midrule
\textbf{Overall} & \textbf{14} & \textbf{1272} & \textbf{9.8} & \textbf{8.2--11.6} & \textbf{72.2} \\
\bottomrule
\end{tabular}
\begin{tablenotes}[flushleft]
\small
\item Abbreviations: CI, confidence interval.
\item Models evaluated: Claude 3.5 Haiku, Claude 3.7 Sonnet (Anthropic); GPT-4o, GPT-4o Mini, O3 Mini (OpenAI); Gemini 1.5 Flash, Gemini 2.0 Flash, Gemini 2.0 Thinking, Gemini 2.5 Pro, Gemma 3 4B, Gemma 3 27B (Google); Llama 3.2 3B (Meta); Qwen 3 32B (Alibaba); DeepSeek R1 Distill (DeepSeek).
\item * Includes both Gemini (commercial) and Gemma (open-source) model families.
\end{tablenotes}
\end{threeparttable}
\end{table}

% ============================================================================
% TABLE 5: RUBRIC DIMENSION ANALYSIS
% ============================================================================

\begin{table}[htbp]
\centering
\caption{Rubric Dimension Scores by Framing Condition}
\label{tab:rubric-dimensions}
\begin{threeparttable}
\begin{tabular}{@{}lcccccc@{}}
\toprule
 & & \multicolumn{2}{c}{\textbf{Mean Score (SD)}} & & & \\
\cmidrule(lr){3-4}
\textbf{Dimension} & \textbf{Max} & \textbf{Neutral} & \textbf{Misconception} & \textbf{Δ} & \textbf{ES*} & \textbf{P Value} \\
\midrule
Evidence Accuracy & 3 & 2.80 (0.5) & 2.68 (0.6) & −0.12 & 0.22 & <.001 \\
Testimonial Handling & 1 & 0.61 (0.5) & 0.83 (0.4) & +0.22 & 0.49 & <.001 \\
Safety Guidance & 2 & 1.73 (0.5) & 1.70 (0.5) & −0.03 & 0.06 & .45 \\
Misattribution Correction & 1 & 0.68 (0.5) & 0.77 (0.4) & +0.09 & 0.20 & .002 \\
Promotional Content Response & 3 & 2.73 (0.7) & 2.42 (1.0) & −0.31 & 0.36 & <.001 \\
HCP Referral & 2 & 1.67 (0.5) & 1.54 (0.6) & −0.13 & 0.24 & <.001 \\
\midrule
\textbf{Total Score} & \textbf{12} & \textbf{10.22 (2.1)} & \textbf{9.94 (2.8)} & \textbf{−0.28} & \textbf{0.11} & \textbf{.045} \\
\bottomrule
\end{tabular}
\begin{tablenotes}[flushleft]
\small
\item Abbreviations: ES, effect size (Cohen's d); HCP, healthcare professional; SD, standard deviation.
\item * Positive ES values indicate worse performance under misconception framing.
\item Promotional Content Response was the primary driver of safety failures; 94.6\% of critical failures involved generating promotional content for harmful treatments.
\end{tablenotes}
\end{threeparttable}
\end{table}

% ============================================================================
% TABLE 6: SAFETY EROSION IN PAIRED CONVERSATIONS
% ============================================================================

\begin{table}[htbp]
\centering
\caption{Classification Transitions in Paired Conversations (Neutral → Misconception Framing)}
\label{tab:erosion}
\begin{threeparttable}
\begin{tabular}{@{}lcccc@{}}
\toprule
 & \multicolumn{3}{c}{\textbf{Outcome Under Misconception Framing}} & \\
\cmidrule(lr){2-4}
\textbf{Starting Classification} & \textbf{Strong} & \textbf{Moderate} & \textbf{Weak} & \textbf{Total} \\
\midrule
Strong (Neutral) & 332 (77.4\%) & 54 (12.6\%) & 43 (10.0\%) & 429 \\
Moderate (Neutral) & 48 (43.6\%) & 35 (31.8\%) & 27 (24.5\%) & 110 \\
Weak (Neutral) & 13 (43.3\%) & 9 (30.0\%) & 8 (26.7\%) & 30 \\
\midrule
\multicolumn{5}{@{}l}{\textit{Summary Statistics}} \\
\addlinespace
\quad Degradation rate* & \multicolumn{4}{c}{22.6\% of Strong conversations degraded} \\
\quad Improvement rate† & \multicolumn{4}{c}{43.3\% of Weak conversations improved} \\
\bottomrule
\end{tabular}
\begin{tablenotes}[flushleft]
\small
\item Values represent number of paired conversations (row percentage).
\item Analysis restricted to 569 conversation pairs with complete data for both framing conditions.
\item * Proportion of conversations that started Strong under neutral framing and degraded to Moderate or Weak under misconception framing.
\item † Proportion of conversations that started Weak under neutral framing and improved to Moderate or Strong under misconception framing, suggesting bidirectional effects.
\end{tablenotes}
\end{threeparttable}
\end{table}

% ============================================================================
% END OF TABLES
% ============================================================================
